\section{The residual-tuple algebraic chart and its lower locus}
\label{sec:residual-tuple-algebraic-chart}

This section proves the algebraization and point interfaces for the
changed residual tuple.  It does not identify this new functor with the
single-character functor in the original source.  The dimension
estimates involving \texttt{Frame.Big}, the arithmetic counting
estimate, and containment in a prescribed Hecke ideal are not
consequences of the chart theorem.

Retain the data and universal ring \(R\) of AC-G94-space.  Thus
\(\mathcal O\) is a complete discrete valuation ring with finite residue
field \(k\), \(G\) is profinite, and \(R\) is a complete Noetherian local
\(\mathcal O\)-algebra with residue field \(k\).  Its universal datum is a
fixed-basis continuous rank-three representation \(r\), together with
the ordered local characters \(\kappa_{v,i}\).  It represents the
normalized polarized ordinary lifts of the finite residual tuple
\(\boldsymbol\theta=(\theta_0,\theta_1,\theta_2)\).
Universality includes every complete Noetherian local test
\(\mathcal O\)-algebra with residue field \(k\), including rings of
characteristic \(p\) and \(R[[U]]\).

Put \(w=(1,0,-1)\), and give the matrix entry \(r_{ij}(g)\) weight
\(w_i-w_j\).  All local characters have weight zero.  Let \(B\) be the
\(\mathfrak m_R\)-adic closure of the \(\mathcal O\)-algebra generated by
all zero-weight monomials in matrix entries and all
\(\kappa_{v,i}(h)\), and put
\[
 A=B[\,r_{ij}(g):g\in G,\ 0\leq i,j\leq2\,]\subset R,\qquad
 \mathfrak m_B=B\cap\mathfrak m_R,\qquad x_0=\mathfrak m_B A.
 \tag{G95.1}
\]
Products of matrix entries, with the sum of their diagonal conjugation
weights, define this grading.  The construction therefore makes sense
for the two cubic adjacent characters.

\begin{theorem}[AC-G95-algebra: finite algebraization of the tuple space]
\label{AC-G95-algebra}
There are finitely many centered universal coordinates \(X_a\in
\mathfrak m_R\), of weights in \(\{-2,-1,0,1,2\}\), such that
\[
 B=R_0,\qquad A=B[X_a],\qquad
 \mathfrak m_R^{\,3m+2}\cap R_j\subset\mathfrak m_B^{\,m}A
       \quad(m\geq1,\ |j|\leq2).
 \tag{G95.2}
\]
Here \(R_j\) is the closed weight-\(j\) part of formal diagonal
conjugation.  The ring \(B\) is complete Noetherian local with residue
field \(k\), and its maximal ideal is the one in (G95.1).
The ring \(A\) is a finitely generated \(B\)-algebra and is
\(\mathbb Z\)-graded, with \(A_0=B\).
Every universal matrix entry is continuous for the \(x_0\)-adic
topology on \(A\).
\end{theorem}

\begin{proof}
\emph{1. Finite homogeneous coordinates from universality.}
Choose the prime-to-\(p\) Teichmuller lift in \(\mathcal O^\times\)
of every nonzero residue value.  Take as coordinates all off-diagonal
entries, all diagonal entries minus the lift of their residual
values, and all local character values minus the lift of their
residual values.  They lie in \(\mathfrak m_R\).
Off-diagonal coordinates have weight \(w_i-w_j\); the centered
diagonal and character coordinates have weight zero.

Their images span
\(\mathfrak m_R/(\mathfrak m_R^2+\mathfrak m_{\mathcal O}R)\).
Indeed a linear functional on this relative cotangent space is the
difference of an \(\mathcal O\)-algebra map to
\(k[\epsilon]/\epsilon^2\) and the residue map.
If it vanishes on all displayed coordinates, the two maps give
the same representation and the same ordered local characters.
Universality says the maps coincide, so the functional is zero.
The cotangent space is finite dimensional because \(R\) is
Noetherian.  Choose a finite spanning subset \(X_a\).

The continuous homomorphism
\[
 \varphi:S=\mathcal O[[Z_1,\ldots,Z_t]]\longrightarrow R,\qquad
 Z_a\longmapsto X_a
 \tag{G95.3}
\]
is surjective.  Directly, the spanning property gives
\(\mathfrak m_R=(\varpi,X_a)R+\mathfrak m_R^2\), hence
\(\mathfrak m_R=(\varpi,X_a)R\) by Nakayama.
Approximate an element successively modulo \(\mathfrak m_R^n\),
using monomials in these generators and coefficients lifting \(k\)
from \(\mathcal O\).  The corrections form a convergent formal
series in (G95.3).

For \(t=1+U\), diagonal conjugation by
\[
 D(t)=\operatorname{diag}(t,1,t^{-1})
 \tag{G95.4}
\]
preserves the representation law, the transpose-antidiagonal
polarization, the normalization at the specified \(g_0\), and the
ordered ordinary matrix identity.  It leaves the local characters
unchanged.  The test ring \(R[[U]]\) is allowed by universality.
Consequently there is a unique homomorphism \(R\to R[[U]]\)
sending each \(X_a\) to \((1+U)^{\mathrm{wt}(X_a)}X_a\).

The kernel of (G95.3) is closed and homogeneous in the formal
sense.  Here is a proof valid even when \(R\) has \(p\)-torsion.
Split a formal series \(f\in S\) into its weight parts \(f_j\).
If \(\varphi(f)=0\), the preceding homomorphism gives
\(\sum_j(1+U)^j\varphi(f_j)=0\).
Modulo \(\mathfrak m_R^n\) only finitely many monomials survive,
so this is a Laurent polynomial in \(t=1+U\).
The substitution
\[
 (R/\mathfrak m_R^n)[t,t^{-1}]
       \longrightarrow(R/\mathfrak m_R^n)[[U]]
 \tag{G95.5}
\]
is injective: multiply by a power of \(t\) and use the
invertible polynomial substitution \(t=1+U\).
Each weight coefficient is zero modulo \(\mathfrak m_R^n\).
Letting \(n\) grow gives \(\varphi(f_j)=0\) for every \(j\).
Thus weight projections descend to \(R\).
They show that a weight-\(j\) element has a weight-\(j\)
preimage in \(S\), also when the element lies in
\(\mathfrak m_R^n\): start with a preimage in
\((\varpi,Z)^n\) and take its weight part.
No division by a weight, a difference of weights, or \(p\)
is involved.

\emph{2. The elementary weight-zero monoid.}
Besides variables of weight zero, the minimal nonempty
zero-weight monomials have weight patterns
\[
 (1,-1),\quad(2,-2),\quad(1,1,-2),\quad(2,-1,-1).
 \tag{G95.6}
\]
This list is exhaustive.  Remove opposite equal weights
if both occur.  If both signs remain, the only possibilities
are \(1,-2\) or \(2,-1\); balance then removes a displayed
triple.  Iteration decomposes every zero-weight monomial.
All indecomposable factors have degree at most three.

There are finitely many such monomials in our variable set.
The zero-weight subring \(S_0\) is the image of the formal
power-series ring over \(\mathcal O\) in these monomials and
the weight-zero variables.
To pass from polynomials to formal series, decompose each
monomial by a fixed rule.  Degree \(d\) requires at least
\(d/3\) factors, so regrouping respects the complete topology.
The same argument includes powers of \(\varpi\), of weight
zero and degree one.  Thus \(S_0\) is complete Noetherian local.

Its image \(R_0\) is closed in \(R\).
A complete Noetherian local ring with finite residue field
is compact, since all quotients by powers of its maximal
ideal are finite.  The image of compact \(S_0\) in Hausdorff
\(R\) is closed.  It follows that \(R_0\) is a complete
Noetherian local quotient of \(S_0\).

Every generator defining \(B\) has weight zero, so
\(B\subset R_0\).  Conversely each centered \(X_a\) is an
actual coordinate or a coordinate minus an \(\mathcal O\)-constant.
Expanding a zero-weight monomial in them puts it in the
algebra defining \(B\).  The closed-image assertion gives
\(R_0\subset B\).  Hence \(B=R_0\).
Its residue map is onto \(k\), since it contains \(\mathcal O\);
its maximal ideal is \(B\cap\mathfrak m_R\).

\emph{3. Bounded remainders and continuity.}
A monomial of weight \(j\), \(|j|\leq2\), from which no
nonempty zero-weight factor can be removed has degree at
most two.  For one sign this is immediate.
For two signs, after excluding opposite equal weights,
only \(1,-2\) or \(2,-1\) remain.
Avoiding a zero-weight triple allows at most one variable
of the lighter opposite weight, and \(|j|\leq2\) gives
degree at most two.
Thus a weight-\(j\) formal series is a finite sum of these
remainder monomials times elements of \(S_0\).

Applying this to the universal entries gives
\(A\subset B[X_a]\); the reverse inclusion follows from
the choice of \(X_a\).
Homogeneity of the relation ideal makes this an algebraic
\(\mathbb Z\)-grading, with degree-zero part \(B\).

For an original monomial of \((\varpi,Z)\)-degree at least
\(3m+2\), removal leaves at most two variables.
Every removed zero-weight factor has degree at most three.
At least \(m\) nonempty zero-weight factors are removed;
their images lie in \(\mathfrak m_B\).
Regrouping by the finitely many remainder monomials proves
(G95.2), also for infinite series.
The homogeneous-preimage observation in Step 1 makes this
an assertion about \(\mathfrak m_R^{3m+2}\cap R_j\).

The original entries are \(\mathfrak m_R\)-adically continuous.
Their differences have the same weight, so (G95.2) proves
\(x_0\)-adic continuity on \(A\).
The same argument applies to local characters, of weight zero.
\end{proof}

\begin{theorem}[AC-G95-chart: the actual lower-triangular point map]
\label{AC-G95-chart}
For a nonzero \(k\)-algebra \(Q\), call an
\(\mathcal O\)-algebra homomorphism \(\xi:A\to Q\) a tuple
\(\mathrm{FPt}\) if \(x_0\subset\ker\xi\);
\(\mathcal O\) acts on \(Q\) through \(k\).
Every such map has fixed diagonal \(\boldsymbol\theta\),
fixed local characters
\(\theta_i\circ\operatorname{loc}_v\), and locally constant
matrix entries.

There is a natural bijection between the tuple
\(\mathrm{FPt}\) maps killing all upper-triangular entries
and the locally constant normalized polarized ordinary
lower-triangular representations over \(Q\) with that
diagonal and those local characters.
Every such representation has a unique actual map \(A\to Q\).
Reading its entries \((1,0)\) and \((2,0)\) gives the
changed-tuple \(\mathrm{At}\) interface.
\end{theorem}

\begin{proof}
For \(b\in B\), choose \(c\in\mathcal O\) whose residue is
the residue of \(b\) in \(R\).  Then \(b-c\in\mathfrak m_B\).
Any \(\mathrm{FPt}\) map sends \(b\) to that specified scalar
in \(Q\).  Diagonal and local character values therefore
have their prescribed residues.
For matrix entries, continuity modulo \(\mathfrak m_R^5\)
and (G95.2) with \(m=1\) prove local constancy after \(\xi\).
All representation, polarization, normalization and ordinary
identities survive every ring homomorphism.

We construct the converse without applying the formal
universal property directly to the generally nonlocal \(Q\).
Let \(r_Q\) be the given locally constant representation.
Compactness gives a finite image.  Intersect its open normal
kernel with its star image and the kernels of the residual
characters.  All data then factor through a finite group
\(\overline G\) with its induced star operation.

Form a polynomial \(k\)-algebra with variables for the
three strictly lower entries of a matrix for each element
of \(\overline G\).  Fix its diagonal to
\(\boldsymbol\theta\) and impose the representation law,
polarization, normalization and ordered ordinary identities.
There are finitely many relations: each of finitely many
local maps has finite image in \(\overline G\).
Let \(D\) be the resulting algebra.
Assign degree one to the adjacent lower entries and
degree two to the long-root entry.
All relations are homogeneous, as follows entry by entry
from matrix multiplication, or from diagonal conjugation.
The diagonal representation gives an augmentation, so
\[
 D=\bigoplus_{d\geq0}D_d,\qquad D_0=k .
 \tag{G95.7}
\]
The universal matrices over \(D\) give a representation
of \(\overline G\), and evaluation gives \(D\to Q\).

Complete \(D\) at \(D_+=\bigoplus_{d>0}D_d\).
The ring \(\widehat D\) is complete Noetherian local
with residue field \(k\), and \(D\to\widehat D\) is injective:
a polynomial has bounded graded degree, whereas a
nonzero homogeneous term in \(D_+^n\) has degree at least \(n\).
The finite-group representation over \(\widehat D\)
is a valid continuous tuple deformation.
Universality supplies
\[
 R\longrightarrow\widehat D .
 \tag{G95.8}
\]
This map is local and continuous: composition with the
residue map onto \(k\) has the unique maximal ideal of \(R\)
as kernel.

A zero-weight monomial in lower-triangular entries is
either a product of diagonal entries or zero; every
nontrivial strictly lower monomial has negative weight.
Local characters are constants.
The defining algebra of \(B\) therefore maps into \(k\).
Continuity and closedness of \(k\subset\widehat D\)
put the entire ring \(B\) in \(k\).
All matrix entries map into \(D\), not just its completion.
Thus (G95.8) restricts to \(A\to D\), and composition with
\(D\to Q\) is the required \(\xi\).
Its restriction to \(B\) is the residue map, so it kills \(x_0\).

A tuple \(\mathrm{FPt}\) map is already determined on
\(B\) by residues, and \(A\) is generated by \(B\)
and all matrix entries.
Agreement of the representations gives uniqueness.
This proves the bijection, its compatibility with maps
of \(k\)-algebras, and the \(\mathrm{At}\) assertion.
\end{proof}

\begin{corollary}[AC-G95-lower-locus: connectedness and dense adjacent points]
\label{AC-G95-lower-locus}
In the actual cubic tuple situation of AC-G90, let
\(\mathfrak t_-\), \(\mathfrak t_+\) and \(\mathfrak e\)
be the ideals of \(A\) generated respectively by the
strictly lower entries, strictly upper entries, and
entries of odd weight.  Put
\[
 \mathcal B=\{\mathfrak p\in\operatorname{Spec}A:
       x_0\subset\mathfrak p,\quad
       \mathfrak t_-\not\subset\mathfrak p\}.
 \tag{G95.9}
\]
Suppose there is one normalized strict polarized filling
with a nonzero adjacent class and one nonzero strict
central \(c\)-anti class, both over \(k\).
Then \(\mathcal B\) is nonempty and connected, every
point kills \(\mathfrak t_+\), and
\[
 \mathcal B\subset
 \overline{\{\mathfrak p\in\mathcal B:
                      \mathfrak e\not\subset\mathfrak p\}} .
 \tag{G95.10}
\]
The two class hypotheses are supplied on the sufficiently
high hosts of AC-G93-fill.
The universal-space hypothesis remains exactly that of
AC-G94-space; its arithmetic group-finiteness input must
be verified, rather than inferred from this corollary.
\end{corollary}

\begin{proof}
Let \(\mathfrak p\in\mathcal B\) and
\(K=\operatorname{Frac}(A/\mathfrak p)\).
A lower entry of some negative weight \(-d\),
\(d\in\{1,2\}\), is nonzero.
If an upper entry of positive weight \(e\in\{1,2\}\)
were nonzero, the product of the \(e\)-th power of
the lower entry and the \(d\)-th power of the upper
entry would be a nonzero zero-weight monomial in \(K\).
In \(A\) it belongs to \(B\cap\mathfrak m_R=\mathfrak m_B\),
so its image is zero, a contradiction.
Thus all upper entries vanish.
The preceding theorem identifies this field point with
an actual locally constant lower representation.

Use the coordinates of AC-G90:
\(r=U(x,y,z)\operatorname{diag}(a,1,d)\).
For such a representation and a strict central
\(c\)-anti cocycle \(b\), the polynomial family is
\[
 x_T=Tx,\qquad y_T=Ty,\qquad z_T=T^2z+(1-T)b .
 \tag{G95.11}
\]
AC-G90 proves every representation, local, polarization
and normalization identity for this family.
AC-G95-chart makes it an actual map \(A\to K[T]\).
If \(x\ne0\) and \(b\ne0\), the lower-entry ideal in
\(K[T]\) is the whole ring: one entry is a nonzero
scalar times \(T\), and a central entry has nonzero
constant term.
Thus the entire image of \(\operatorname{Spec}K[T]\)
lies in \(\mathcal B\).
At \(T=0\) it is the central representation \(b\);
at \(T=1\) it is the chosen lower representation.
For \(T\ne0\) an adjacent entry is nonzero.

Fix a nonzero central point \(b_0\) and a filled adjacent
point \(r_*\) over \(k\).
Every field point with a nonzero adjacent entry is
joined to \(b_0\) by (G95.11).
If its adjacent entries vanish, its nonzero lower
entry is central; use it as \(b\) and use the base
change of \(r_*\) for the other endpoint.
This joins it to \(r_*\), which is joined to \(b_0\).
Each image is connected because
\(\operatorname{Spec}K[T]\) is connected.
Base change of a fixed \(k\)-point has the same kernel
as that \(k\)-point.
These connected chains prove connectedness of the
entire subspace \(\mathcal B\) and its nonemptiness.

For a central point, the last family has adjacent
entries nonzero on the dense open set \(T\ne0\).
Continuity on spectra puts its \(T=0\) point in the
closure of that locus.
Points already in the adjacent locus need no
specialization.  This proves (G95.10).
No irreducibility of the full mixed-cup zero locus
is asserted or used.
\end{proof}

\paragraph{Exact source interface and geometric scope.}
The original definitions are in
\path{LeAdicClosure.lean}: \texttt{NRep.IsBar},
\texttt{Space}, \texttt{Space.B}, \texttt{Space.A},
\texttt{NRep.mon} and the restricted universal representation;
\path{StrSMul.lean}: \texttt{Space.FPt} and \texttt{Space.At};
and the theorem interfaces in
\path{SmoothCharacterExtendsToOpenSubgroup.lean},
\path{SmoothHomWeightSumNeg.lean} and
\path{ExistsNonzeroLowCocycle.lean}.
All are from the pinned source commit
\path{7cbf7a0c32fca1badf05931879a37c3f8a2b03b4}.
This proof constructs their changed-tuple counterparts
directly.  It does not apply a theorem whose original
\texttt{IsBar} still demands powers of one character.
The formal torus and finite graded algebra arguments
are proved above; no source dimension theorem is
imported by analogy.

Geometrically \(R\) is the formal deformation space
at the diagonal representation, \(B\) is its
completed weight-zero coordinate ring, and
\(\operatorname{Spec}A\) is the algebraic chart
generated by representation entries over that ring.
The maps in (G95.11) are actual morphisms from
affine lines in the closed invariant fiber.
This uses the established formal-scheme and
representation-functor dictionary, without claiming
a new EGA or archival scan reading.
The result supplies the chart and lower-locus topology.
It supplies no adjacent class-sector upper bound,
no \texttt{Frame.Big} dimension inequality, no
original counting budget, and no containment of a
prescribed curve in the native Hecke support.
